% Source: https://artofproblemsolving.com/community/c3h2534090_interesting_question
% Stable identity: aops-topic-2534090-q01
\textbf{(a)} \emph{Existence.} Put
\[P(x)=y_1\dfrac{(x-x_2)(x-x_3)}{(x_1-x_2)(x_1-x_3)}+y_2\dfrac{(x-x_1)(x-x_3)}{(x_2-x_1)(x_2-x_3)}+y_3\dfrac{(x-x_1)(x-x_2)}{(x_3-x_1)(x_3-x_2)}.\]
The denominators are nonzero because the $x_i$ are different. At $x=x_1$ the second and third fractions vanish and the first equals $1$, so $P(x_1)=y_1$; similarly $P(x_2)=y_2$ and $P(x_3)=y_3$. So $P(x)=ax^2+bx+c$ is a polynomial of degree at most $2$ through the three points (Lagrange interpolation). If its coefficient $a$ were $0$, its graph would be a line containing all three points, contradicting non-collinearity. So $a\ne0$ and $y=P(x)$ is a parabola through the points.

\emph{Uniqueness.} If $Q(x)$ is another polynomial of degree at most $2$ with $Q(x_i)=y_i$, then $P-Q$ has degree at most $2$ and three different roots $x_1,x_2,x_3$. A nonzero polynomial of degree at most $2$ has at most $2$ roots, so $P-Q$ is the zero polynomial and $Q=P$. $\blacksquare$

(The same argument shows that $n+1$ points with different $x$-coordinates determine exactly one polynomial of degree at most $n$.)

\textbf{(b)} Substituting the points into $y=ax^2+bx+c$:
\[a-b+c=2,\qquad a+b+c=0,\qquad9a+3b+c=4.\]
Subtracting the first two gives $b=-1$, so $a+c=1$; then $9a-3+c=4$ gives $8a=6$. Hence $a=\dfrac34$, $c=\dfrac14$, and the parabola is
\[\boxed{y=\dfrac34x^2-x+\dfrac14}.\]

\begin{asy}
usepackage("amsmath");
size(4.5cm);
defaultpen(fontsize(10pt));
import graph;
real f(real x){return 0.75*x^2-x+0.25;}
draw((-2.2,0)--(4.2,0),Arrow(4));draw((0,-0.6)--(0,6.2),Arrow(4));
draw(graph(f,-1.6,3.4),linewidth(0.8));
pair P1=(-1,2),P2=(1,0),P3=(3,4);
dot(P1);dot(P2);dot(P3);
label("$(x_1,y_1)$",P1,W);label("$(x_2,y_2)$",P2,SE);label("$(x_3,y_3)$",P3,E);
label("$x$",(4.2,0),S);label("$y$",(0,6.2),W);
\end{asy}
